\documentclass[12pt,a4paper]{article}

% ---------- پکیج‌ها (همه قبل از xepersian) ----------
\usepackage{fontspec}
\usepackage{amsmath, amssymb, amstext}
\usepackage{graphicx}
\usepackage{geometry}
\usepackage{xcolor}
\usepackage{fancyhdr}

% ---------- تنظیم حاشیه ----------
\geometry{margin=1in}

% ---------- تنظیم fancyhdr ----------
\pagestyle{fancy}
\setlength{\headheight}{15pt}

% هدرها
\lhead{کوییز درس فیزیک \lr{II}}
\chead{دانشکده فیزیک دانشگاه صنعتی خواجه نصیرالدین طوسی}
\rhead{\today}

% ---------- بسته فارسی و فونت ----------
\usepackage{xepersian}
\settextfont{B Nazanin}

\begin{document}
	
	% ---------- جدول اطلاعات ----------
	\begin{center}
		\begin{tabular}{|p{7cm}|p{7cm}|}
			\hline
			\textbf{نام و نام خانوادگی:} & \textbf{شماره دانشجویی:} \\
			\hline
			\vspace{0.5cm} & \vspace{0.5cm} \\
			\hline
		\end{tabular}
	\end{center}
	
	\vspace{10pt}
	
	% ---------- سوالات ----------
	\section*{سوالات:}
	
	\begin{enumerate}
		
		% ---------- سوال ۱ ----------
		\item
		یک حلقه نازک و مسطح به شعاع $R$ و بار کل $Q$ را در نظر بگیرید.
		\begin{enumerate}
			\item[(الف)] پتانسیل الکتریکی را در نقطه‌ای واقع بر محور حلقه و در فاصله $z$ از مرکز آن به دست آورید.
			\item[(ب)] میدان الکتریکی را در همان نقطه محاسبه کنید.
		\end{enumerate}
		
		\vspace{6pt}
		
		% ---------- سوال ۲ ----------
		\item
		یک کره نارسانای توپر به شعاع $a$ را در نظر بگیرید که بار مثبت کل $Q$
		به‌طور یکنواخت در حجم آن توزیع شده است.
		\begin{enumerate}
			\item[(الف)] پتانسیل الکتریکی را در ناحیه داخل کره $(r<a)$ به دست آورید.
			\item[(ب)] پتانسیل الکتریکی را در ناحیه بیرون کره $(r>a)$ به دست آورید.
		\end{enumerate}
		
	\end{enumerate}
	
	\vspace{12pt}
	\begin{center}
		\textbf{موفق باشید }
	\end{center}
	
\end{document}
